% !TEX root = ../thesis.tex
\chapter{Phương pháp đề xuất}

Chương này trình bày chi tiết thiết kế và hiện thực hóa kiến trúc pipeline HLK-ViT5.
Mục~3.1 tổng quan luồng xử lý toàn hệ thống;
Mục~3.2 mô tả bộ dữ liệu thực nghiệm 15.620 mẫu từ \textit{tienphong.vn};
Mục~3.3 chi tiết hóa 5 module cải tiến cốt lõi;
Mục~3.4 trình bày thiết lập siêu tham số huấn luyện; và Mục~3.5 tổng kết chương.

\section{Tổng quan kiến trúc Pipeline HLK-ViT5}

Pipeline HLK-ViT5 được xây dựng nhằm giải quyết bài toán tóm tắt văn bản tin tức ngữ cảnh dài thông qua quy trình xử lý 13 bước được tự động hóa hoàn toàn. Kiến trúc tổng thể của hệ thống được thể hiện trong Sơ đồ bên dưới.

\begin{figure}[H]
  \centering
  \begin{tikzpicture}[
    auto, >=stealth',
    block/.style={rectangle, draw, fill=blue!10, text width=8em, text centered, rounded corners, inner sep=4pt, font=\small},
    cloud/.style={draw, ellipse, fill=orange!15, inner sep=3pt, font=\small},
    line/.style={draw, ->, thick}]

    % Column 1: Preprocessing
    \node [cloud] (input) {Văn bản gốc};
    \node [block, below=0.5cm of input] (norm) {1. Chuẩn hóa \& Xử lý khoảng trắng};
    \node [block, below=0.5cm of norm] (split) {2. Tách câu chuẩn};
    \node [block, below=0.5cm of split] (chunk) {3. Tạo Chunk $\le 256$ token (Overlap 1 câu)};

    % Column 2: Context Filtering
    \node [block, right=1.0cm of norm] (score) {4. Chấm điểm: Coverage, Salience, Position};
    \node [block, below=0.5cm of score] (mmr) {5. Lọc MMR giảm trùng lặp};
    \node [block, below=0.5cm of mmr] (select) {6. Chọn Max 4 Chunk \& Khôi phục vị trí};

    % Column 3: Generation & Reranking
    \node [block, right=1.0cm of score] (prompt) {7. Tạo Input (Keywords + Context $\le 1024$)};
    \node [block, below=0.5cm of prompt, fill=green!15] (vit5) {8. ViT5-base Model};
    \node [block, below=0.5cm of vit5] (rerank) {9. Sinh Candidate \& Factuality Reranker};
    \node [cloud, below=0.5cm of rerank, fill=red!15] (output) {Bản tóm tắt cuối cùng};

    % Paths & Arrows
    \path [line] (input) -- (norm);
    \path [line] (norm) -- (split);
    \path [line] (split) -- (chunk);
    \path [line] (chunk.east) -- ++(0.4,0) |- (score.west);
    \path [line] (score) -- (mmr);
    \path [line] (mmr) -- (select);
    \path [line] (select.east) -- ++(0.4,0) |- (prompt.west);
    \path [line] (prompt) -- (vit5);
    \path [line] (vit5) -- (rerank);
    \path [line] (rerank) -- (output);

  \end{tikzpicture}
  \caption{Sơ đồ luồng xử lý 13 bước trong kiến trúc pipeline HLK-ViT5.}
  \label{fig:hlk_vit5_pipeline}
\end{figure}

Luồng xử lý đi từ văn bản bài báo gốc, trải qua công đoạn tiền xử lý, phân đoạn nén ngữ cảnh thông minh bằng MMR, ghép chuỗi điều kiện hóa từ khóa, đưa qua mô hình ViT5-base và cuối cùng lọc qua bộ xếp hạng lại tính xác thực để xuất bản tóm tắt tối ưu.

\section{Bộ dữ liệu thực nghiệm}

Nghiên cứu xây dựng và sử dụng bộ dữ liệu corpus tin tức tiếng Việt được thu thập trực tiếp từ trang báo điện tử \textit{tienphong.vn}. Bộ dữ liệu bao gồm các đặc điểm chính:

\begin{itemize}
  \item \textbf{Quy mô tập huấn luyện:} 15.620 cặp văn bản gốc và bản tóm tắt tiêu chuẩn.
  \item \textbf{Đặc trưng ngữ cảnh:} Độ dài văn bản gốc trung bình từ 800 đến 1.500 từ, bao phủ nhiều thể loại báo chí khác nhau như Thời sự, Kinh tế, Pháp luật, Văn hóa và Xã hội.
  \item \textbf{Chuẩn hóa dữ liệu:} Tất cả văn bản được xử lý khoảng trắng thừa, sửa lỗi mã hóa Unicode dựng sẵn và tách từ tiếng Việt.
\end{itemize}

\section{Chi tiết các Module cải tiến cốt lõi}

\subsection{Sentence-aware Chunking và Overlapping}

Bài báo gốc sau khi chuẩn hóa được tách thành danh sách các câu hoàn chỉnh $S = (s_1, s_2, \ldots, s_K)$. Quá trình phân đoạn được thực hiện bằng cách gom các câu liên tiếp sao cho tổng số token của mỗi chunk $C_j$ không vượt quá 256 tokens:

\[
|C_j| \;\le\; 256 \text{ tokens}
\]

Nhằm tránh làm gãy đứt ngữ cảnh giao thoa giữa hai đoạn kề nhau, chunk $C_{j+1}$ luôn giữ lại câu cuối cùng của chunk $C_j$ làm câu mở đầu (overlap 1 câu).

\subsection{Salience Scoring và Lọc ngữ cảnh bằng MMR}

Mỗi chunk $C_j$ được tính điểm tổng hợp dựa trên 3 thành phần:

\[
\text{Score}(C_j) \;=\; \alpha \cdot \text{KeywordCoverage}(C_j) + \beta \cdot \text{Salience}(C_j) + \gamma \cdot \text{PositionScore}(C_j)
\]

\noindent trong đó:
\begin{itemize}
  \item $\text{KeywordCoverage}(C_j)$ đo tỷ lệ từ khóa chủ đề xuất hiện trong chunk.
  \item $\text{Salience}(C_j)$ tính độ nổi bật thông tin dựa trên điểm TF-IDF trung bình của các thuật ngữ.
  \item $\text{PositionScore}(C_j)$ thưởng điểm cho các chunk nằm ở đoạn đầu và đoạn kết của bài báo.
\end{itemize}

Sau khi tính điểm, thuật toán MMR chọn lọc ra tối đa 4 chunk có điểm số cao nhất nhưng không trùng lặp ngữ nghĩa với các chunk đã chọn. Bốn chunk này sau đó được khôi phục về đúng vị trí xuất hiện ban đầu trong bài báo để đảm bảo tính tuyến tính của mạch văn.

\subsection{Input Formatting và Keyword Conditioning}

Bốn chunk chọn lọc được nối lại thành ngữ cảnh đại diện $\tilde{X}$. Danh sách từ khóa quan trọng $K = \{k_1, k_2, \ldots, k_m\}$ được trích xuất và đặt ở vị trí đầu chuỗi prompt:

\begin{center}
\texttt{tóm tắt: Từ khóa: } $k_1, k_2, \ldots, k_m$ \texttt{ | Ngữ cảnh: } $\tilde{X}$
\end{center}

Cấu trúc input này mở rộng cửa sổ ngữ cảnh hiệu quả của ViT5 lên đến 1024 tokens, cho phép mô hình tiếp nhận thông tin bao quát toàn văn bài báo.

\subsection{Fine-tuning mô hình ViT5-base và phát hành Checkpoint}

Mô hình \texttt{VietAI/ViT5-base} được tinh chỉnh có giám sát trên tập 15.620 mẫu với hàm mất mát Cross-Entropy chuẩn cho bài toán sinh chuỗi:

\[
\mathcal{L}_{\text{CE}}(\theta) \;=\; -\sum_{t=1}^{T} \log P_\theta(y_t \mid y_{<t}, \tilde{X}, K)
\]

Trọng số mô hình sau khi huấn luyện tinh chỉnh đã được đóng gói và phát hành công khai trên kho lưu trữ Hugging Face Hub tại địa chỉ: \url{https://huggingface.co/tmanh217/hlk-vit5}.

\subsection{Candidate Generation và Heuristic Factuality Reranker}

Trong giai đoạn suy luận, thay vì chỉ sinh 1 bản tóm tắt bằng giải thuật Beam Search thông thường, ViT5-base được cấu hình sinh ra $N = 5$ ứng viên tóm tắt $\{Y_1, Y_2, \ldots, Y_N\}$.

Bộ xếp hạng lại Heuristic Factuality Reranker chấm điểm từng ứng viên $Y_i$ dựa trên hai tiêu chí:

\begin{itemize}
  \item \textbf{Entity Consistency Score:} Phạt nặng các ứng viên xuất hiện tập thực thể $E(Y_i)$ không nằm trong tập thực thể bài gốc $E(X)$.
  \item \textbf{Number Accuracy Score:} So sánh tập các con số $N(Y_i)$ với $N(X)$, loại bỏ các bản tóm tắt làm sai lệch số liệu.
\end{itemize}

Bản tóm tắt $Y^{\star}$ có điểm Reranker cao nhất được chọn làm kết quả xuất ra cuối cùng.

\section{Thiết lập thực nghiệm và Siêu tham số}

Bảng~\ref{tab:training_hyperparams} tóm tắt các siêu tham số chính được sử dụng trong quá trình huấn luyện HLK-ViT5.

\begin{table}[htbp]
\centering
\caption{Cấu hình siêu tham số huấn luyện mô hình HLK-ViT5}
\label{tab:training_hyperparams}
\begin{tabularx}{\textwidth}{l X}
\hline
\textbf{Siêu tham số} & \textbf{Giá trị thiết lập} \\ \hline
Mô hình nền & \texttt{VietAI/ViT5-base} \\
Hugging Face Checkpoint & \url{https://huggingface.co/tmanh217/hlk-vit5} \\
Độ dài đầu vào tối đa & 1024 tokens \\
Độ dài đầu ra tối đa & 256 tokens \\
Tốc độ học (Learning Rate) & $5 \times 10^{-5}$ (với AdamW optimizer) \\
Kích thước Batch (Batch Size) & 16 (với gradient accumulation) \\
Số lượng Epochs & 5 epochs \\
Số lượng Candidate Rerank ($N$) & 5 candidates \\
Hệ số MMR ($\lambda$) & 0,7 \\ \hline
\end{tabularx}
\end{table}

\section{Tổng kết chương}

Chương này đã chi tiết hóa kiến trúc 13 bước của HLK-ViT5, kết hợp giữa chia đoạn bảo toàn câu, nén ngữ cảnh MMR, Keyword Conditioning và Heuristic Factuality Reranker. Toàn bộ thiết lập này được huấn luyện trên 15.620 mẫu dữ liệu \textit{tienphong.vn}. Chương~4 sẽ báo cáo kết quả thực nghiệm định lượng.
